\documentclass[conference]{IEEEtran}

\usepackage[T1]{fontenc}
\usepackage[utf8]{inputenc}
\usepackage{cite}
\usepackage{url}
\usepackage{booktabs}
\usepackage{array}
\usepackage{tabularx}
\usepackage{hyperref}
\usepackage{xcolor}

\hypersetup{
  colorlinks=true,
  linkcolor=black,
  citecolor=black,
  urlcolor=black
}

\newcommand{\srs}[1]{\textbf{#1}}
\newcommand{\RoomLink}{RoomLink}

\title{RoomLink: A QR-Integrated University Room Reservation System}

\author{
\IEEEauthorblockN{Ibrahim Utlu}
\IEEEauthorblockA{
Department of Computer Engineering\\
Istanbul Medipol University\\
Istanbul, Turkey}
}

\begin{document}

\maketitle

\begin{abstract}
This paper reports the design, implementation, and deployment of \RoomLink, a
QR-integrated university room reservation system. The system supports student
room discovery, time-slot reservation, QR-based check-in and check-out, break
mode, notifications, and role-based administrative workflows. The report is
structured to remain directly aligned with the Software Requirements
Specification (SRS). Since the SRS is still being finalized, this document uses
a preliminary SRS baseline derived from the current repository documentation,
API contract, database schema, and implemented user interface. Each major
design and implementation decision is mapped to explicit requirement
identifiers, allowing the final SRS identifiers to be substituted without
rewriting the technical narrative. The resulting system uses an ASP.NET Core
Web API, PostgreSQL, a static HTML/CSS/JavaScript frontend, JSON Web Token
authentication, QR code generation and validation, and cloud deployment through
Railway and Vercel.
\end{abstract}

\begin{IEEEkeywords}
room reservation, QR code, software requirements specification, ASP.NET Core,
PostgreSQL, role-based access control, university information system
\end{IEEEkeywords}

\section{Introduction}
Universities often need shared room management for study spaces, presentation
rooms, classrooms, and staff-controlled facilities. Manual reservation
processes can create conflicts, lack reliable attendance tracking, and provide
limited visibility into real-time room usage. \RoomLink\ addresses these
problems by combining a web-based reservation workflow with QR-based physical
presence verification.

The system allows students to browse rooms, reserve available time slots, and
scan room QR codes to check in, start or end a break, and check out. Staff
members can monitor reservations and view live QR codes, while administrators
can manage rooms, users, reservations, and QR rotation. The implementation is
therefore not only a booking interface but also a lightweight operational
control system for university rooms.

This report follows an IEEE-style engineering paper structure and is written
to be consistent with the SRS. Because the SRS is near completion but not yet
final, all requirements are referenced through stable preliminary identifiers.
The final SRS can later adopt these identifiers or map them to its official
numbering.

\section{SRS Alignment Method}
\subsection{Preliminary SRS Baseline}
No standalone final SRS document is currently present in the repository.
Therefore, this report treats the following project artifacts as the
preliminary SRS baseline \cite{roomlink_readme,roomlink_api,roomlink_database,roomlink_backend}:
\begin{itemize}
  \item project overview and role descriptions in the repository README;
  \item API behavior and endpoint definitions in the API endpoint document;
  \item database tables, status values, and constraints in the database
  documentation;
  \item backend service descriptions and reservation rules in the backend
  documentation;
  \item implemented frontend pages for login, rooms, reservation creation,
  reservation viewing, QR scanning, administration, and QR monitoring.
\end{itemize}

\subsection{Requirement Identifier Convention}
The report uses the following SRS identifier convention:
\begin{itemize}
  \item \srs{SRS-FR-x}: functional requirements;
  \item \srs{SRS-NFR-x}: non-functional requirements;
  \item \srs{SRS-DR-x}: data requirements;
  \item \srs{SRS-UC-x}: use-case requirements;
  \item \srs{SRS-SEC-x}: security and access-control requirements.
\end{itemize}

When the SRS is finalized, these identifiers should be checked against the
official SRS numbering. Any mismatch should be resolved in Table
\ref{tab:traceability} before submission. This traceability-first approach is
consistent with requirements engineering guidance in ISO/IEC/IEEE 29148
\cite{ieee29148}.

\section{System Scope and Stakeholders}
\RoomLink\ serves three primary stakeholder groups. Students use the system to
reserve rooms and prove room attendance through QR scans. Staff use the system
to monitor reservations and display live QR codes. Administrators perform
system management tasks, including room CRUD operations, student registration,
reservation monitoring, and QR rotation.

Table \ref{tab:roles} summarizes the role-to-capability model. This model is
used consistently in the user interface, backend authorization policy, and API
documentation.

\begin{table}[htbp]
\caption{Stakeholders and Core Capabilities}
\label{tab:roles}
\centering
\begin{tabularx}{\linewidth}{@{}lX@{}}
\toprule
\textbf{Role} & \textbf{Capabilities} \\
\midrule
Student & Browse rooms, create reservations, view own reservations, scan QR
codes, start and end break mode, receive notifications. \\
Staff & Monitor active reservations, view room QR codes, use the live QR
monitor, access staff-level operational screens. \\
Admin & All staff capabilities plus student registration, room management,
reservation monitoring, and QR rotation. \\
\bottomrule
\end{tabularx}
\end{table}

\section{Functional Requirements}
The implemented system is organized around the functional requirements listed
in Table \ref{tab:functional}. These requirements are intentionally written in
SRS language so that the report remains traceable to the final requirements
document.

\begin{table*}[htbp]
\caption{Functional Requirements Baseline}
\label{tab:functional}
\centering
\begin{tabularx}{\textwidth}{@{}lXl@{}}
\toprule
\textbf{ID} & \textbf{Requirement} & \textbf{Implemented Surface} \\
\midrule
\srs{SRS-FR-01} & The system shall authenticate students, staff, and
administrators using role-aware login flows. & Login page, Auth API \\
\srs{SRS-FR-02} & The system shall allow authorized users to browse and search
available rooms. & Rooms page, Room API \\
\srs{SRS-FR-03} & The system shall allow users to create room reservations
subject to slot, capacity, and status constraints. & Reserve page,
Reservation API \\
\srs{SRS-FR-04} & The system shall enforce standard two-hour reservation slots
for standard rooms between 08:00 and 22:00. & Room slots, Reservation API \\
\srs{SRS-FR-05} & The system shall allow demo rooms to use flexible start and
end times while still enforcing capacity. & Reserve page, Reservation API \\
\srs{SRS-FR-06} & The system shall allow users to view active reservations
belonging to their account. & My Reservations page \\
\srs{SRS-FR-07} & The system shall allow reservation owners or administrators
to cancel reservations. & Reservation detail/admin pages \\
\srs{SRS-FR-08} & The system shall validate QR scans for check-in, break-out,
break-in, and check-out operations. & Scan page, QR API \\
\srs{SRS-FR-09} & The system shall support rotating dynamic QR codes for rooms.
& QR monitor, QR service \\
\srs{SRS-FR-10} & The system shall notify users about reservation lifecycle
events and warnings. & Notification bar, Notification API \\
\srs{SRS-FR-11} & The system shall allow administrators to manage rooms and
student accounts. & Admin dashboard \\
\srs{SRS-FR-12} & The system shall maintain scan logs for QR-based attendance
actions. & Database, QR controller \\
\bottomrule
\end{tabularx}
\end{table*}

\section{Non-Functional Requirements}
\RoomLink\ also satisfies several non-functional requirements relevant to a
coursework-grade production deployment. The system uses a web architecture that
can be deployed without a frontend build pipeline and keeps backend state in
PostgreSQL. Cloud deployment targets are Railway for the API and database and
Vercel for the static frontend.

\begin{table}[htbp]
\caption{Non-Functional Requirements Baseline}
\label{tab:nfr}
\centering
\begin{tabularx}{\linewidth}{@{}lX@{}}
\toprule
\textbf{ID} & \textbf{Requirement} \\
\midrule
\srs{SRS-NFR-01} & The system shall be accessible through a web browser on
desktop and mobile devices. \\
\srs{SRS-NFR-02} & The system shall use secure token-based authentication for
protected API endpoints. \\
\srs{SRS-NFR-03} & The system shall separate frontend, backend, and database
layers for deployment maintainability. \\
\srs{SRS-NFR-04} & The system shall support cloud deployment with configurable
environment variables. \\
\srs{SRS-NFR-05} & The system shall preserve production data during incremental
schema repair operations. \\
\srs{SRS-NFR-06} & The user interface shall provide responsive layouts for
common mobile and desktop viewport sizes. \\
\bottomrule
\end{tabularx}
\end{table}

\section{Architecture}
\subsection{High-Level Architecture}
The architecture follows a three-layer web application model. The frontend is
a static web client implemented with HTML, CSS, and vanilla JavaScript. It
communicates with the backend through JSON APIs. The backend is an ASP.NET
Core Web API that implements authentication, authorization, reservation logic,
QR logic, notification creation, and schema repair. PostgreSQL stores users,
rooms, reservations, QR code records, scan logs, and notifications.

\begin{table}[htbp]
\caption{Technology Stack}
\label{tab:stack}
\centering
\begin{tabularx}{\linewidth}{@{}lX@{}}
\toprule
\textbf{Layer} & \textbf{Technology} \\
\midrule
Frontend & Static HTML, CSS, JavaScript, html5-qrcode \\
Backend & ASP.NET Core 5 Web API, Entity Framework Core, JWT Bearer \\
Database & PostgreSQL with raw SQL schema and repair service \\
QR & QRCoder for image generation and dynamic token validation \\
Deployment & Railway API and PostgreSQL, Vercel frontend \\
\bottomrule
\end{tabularx}
\end{table}

\subsection{Component Responsibilities}
The authentication component issues role-bearing JWT tokens. The room component
returns room lists, room status, and available slots. The reservation component
creates, updates, cancels, and retrieves reservation records while enforcing
business rules. The QR component validates scan actions and transitions
reservation states. The notification component exposes user-facing messages.
The schema repair component safely aligns older production databases with the
expected schema.

\section{Data Design}
The database design directly supports the SRS requirement set. The main tables
are \texttt{users}, \texttt{rooms}, \texttt{reservations}, \texttt{qr\_codes},
\texttt{scan\_logs}, and \texttt{notifications}. The reservation model stores
room, user, date, start time, end time, status, and QR payload information.
The scan log table records attendance events, supporting auditability for
check-in, break, and check-out flows.

Reservation status values include \texttt{Pending}, \texttt{Confirmed},
\texttt{Active}, \texttt{CheckedIn}, \texttt{OnBreak}, \texttt{CheckedOut},
\texttt{Cancelled}, \texttt{Expired}, and \texttt{NoShow}. These values map to
the reservation lifecycle required by \srs{SRS-FR-03}, \srs{SRS-FR-08}, and
\srs{SRS-FR-10}.

\section{Business Rules}
The most important business rules are enforced on the backend rather than only
in the user interface. Standard rooms use fixed two-hour slots beginning at
08:00, 10:00, 12:00, 14:00, 16:00, 18:00, and 20:00. Live statuses hold room
capacity and prevent conflicting reservations. Standard rooms enforce one live
reservation per user. Demo rooms are explicitly marked with \texttt{IsDemoRoom}
and bypass the fixed-slot and single-active-reservation rules, while still
enforcing room capacity.

The QR workflow validates that the QR value is acceptable, that the caller has
a matching reservation, and that the reservation status and time window permit
the requested scan action. The QR code concept follows the QR Code symbology
standard \cite{iso_qr}. Break mode keeps a reservation slot held while the
student is away and creates notifications when the break lifecycle requires
attention.

\section{Security and Access Control}
The system uses JWT bearer authentication for protected endpoints. Each user
has a role: Student, Staff, or Admin. Role-based access is enforced in the API.
Students cannot retrieve staff/admin QR display endpoints. Staff users can
monitor operational data but cannot perform admin-only room management.
Administrators can register students, manage rooms, rotate QR codes, and view
all reservations. JWT is used according to the token-based authentication model
defined in RFC 7519 \cite{rfc7519}.

Table \ref{tab:security} maps security requirements to implementation
mechanisms.

\begin{table}[htbp]
\caption{Security Requirement Mapping}
\label{tab:security}
\centering
\begin{tabularx}{\linewidth}{@{}lX@{}}
\toprule
\textbf{ID} & \textbf{Implementation} \\
\midrule
\srs{SRS-SEC-01} & Protected endpoints require a valid JWT bearer token. \\
\srs{SRS-SEC-02} & API policies restrict staff/admin-only QR display
endpoints. \\
\srs{SRS-SEC-03} & Reservation detail and user reservation endpoints are
limited to the owner or administrator. \\
\srs{SRS-SEC-04} & Admin-only endpoints protect student registration, room
CRUD operations, and QR rotation. \\
\srs{SRS-SEC-05} & CORS origins are configurable through deployment
environment variables. \\
\bottomrule
\end{tabularx}
\end{table}

\section{Implementation}
\subsection{Frontend}
The frontend is implemented as static pages. The main user-facing pages are
login, rooms, reservation creation, my reservations, reservation details, QR
scanner, admin dashboard, and live QR monitor. The frontend uses shared CSS for
responsive layout, theme selection, dark mode, message boxes, cards, navigation
bars, and reservation grids. QR scanning is handled through the browser camera
using the html5-qrcode library.

\subsection{Backend}
The backend exposes REST-style API endpoints through ASP.NET Core controllers.
Authentication endpoints handle student and staff/admin login. Room endpoints
return room lists, status, and available slots. Reservation endpoints enforce
the slot and capacity rules. QR endpoints validate check-in, break, and
check-out actions. Notification endpoints return unread and historical
messages.

\subsection{Deployment}
The backend runs as a containerized service on Railway and binds to the
platform-provided port. The frontend is deployed as static files on Vercel.
The backend reads production configuration from environment variables,
including database connection strings, JWT secret, frontend CORS origin, QR
rotation interval, check-in grace period, break duration, and local time-zone
offset.

\section{Testing and Validation}
Testing is organized around the SRS use cases. Student tests validate login,
room browsing, slot selection, reservation creation, duplicate reservation
rejection, QR check-in, QR check-out, and notification display. Staff tests
validate QR visibility, live QR monitor access, and restricted access to
admin-only room management. Admin tests validate student registration, room
CRUD operations, reservation monitoring, QR rotation, and all-reservation
visibility.

The deployment smoke test validates backend health, authentication, room list
retrieval, frontend delivery, and CORS behavior. These tests are aligned with
\srs{SRS-UC-01} through \srs{SRS-UC-09}.

\section{SRS Traceability Matrix}
Table \ref{tab:traceability} summarizes how the report and implementation map
to the preliminary SRS. This table is the primary control point for keeping
the IEEE report synchronized with the final SRS.

\begin{table*}[htbp]
\caption{SRS Traceability Matrix}
\label{tab:traceability}
\centering
\begin{tabularx}{\textwidth}{@{}lXlX@{}}
\toprule
\textbf{SRS ID} & \textbf{Requirement Summary} & \textbf{Implementation Evidence} & \textbf{Report Sections} \\
\midrule
\srs{SRS-FR-01} & Role-aware authentication & Login page, Auth controller & III, VIII \\
\srs{SRS-FR-02} & Room browsing and search & Rooms page, Room controller & III, VI, IX \\
\srs{SRS-FR-03} & Reservation creation & Reserve page, Reservation controller & III, VII, IX \\
\srs{SRS-FR-04} & Two-hour standard slots & Available-slots endpoint, validation rules & III, VII \\
\srs{SRS-FR-05} & Demo room flexible slots & Demo room flag and reservation validation & III, VII \\
\srs{SRS-FR-06} & View own reservations & My Reservations page, user reservation endpoint & III, VIII \\
\srs{SRS-FR-07} & Cancel reservations & Reservation detail and admin dashboard & III, VIII \\
\srs{SRS-FR-08} & QR scan workflow & Scan page, QR controller, scan logs & III, VII, VIII \\
\srs{SRS-FR-09} & Dynamic QR rotation & QR service, QR monitor, rotate endpoint & III, VI, VIII \\
\srs{SRS-FR-10} & Notifications and warnings & Notification service and notification bar & III, VIII, IX \\
\srs{SRS-FR-11} & Admin management & Admin dashboard, room/user endpoints & III, VIII \\
\srs{SRS-SEC-01} & Token authentication & JWT bearer auth & VIII \\
\srs{SRS-NFR-04} & Cloud deployment & Railway, PostgreSQL, Vercel & VI, VIII \\
\bottomrule
\end{tabularx}
\end{table*}

\section{Limitations and Future Work}
The SRS is not yet final, so the requirement identifiers in this report should
be treated as a controlled preliminary baseline. When the final SRS is
completed, the traceability matrix must be reviewed line by line. Potential
future improvements include stronger password hashing, formal EF Core
migrations, richer automated test coverage, refined mobile navigation,
calendar-based reservation views, and analytics for room utilization.

\section{Conclusion}
\RoomLink\ implements a QR-integrated university room reservation workflow
with role-aware authentication, reservation constraints, QR-based attendance,
notifications, administrative controls, and cloud deployment. The report is
written to remain synchronized with the SRS through explicit requirement IDs
and a traceability matrix. This structure allows the final SRS to be completed
with minimal report rework while preserving a clear engineering explanation of
the implemented system.

\bibliographystyle{IEEEtran}
\bibliography{references}

\end{document}
